\section{Round-twenty-three reconstruction: relative canonical normalization and the full mixed local coefficient}
\label{sec:r23-b1}

The previous canonical pressure was defined from an unnormalised
$1/N!$ integral and consequently contained the divergent ideal-gas term
$-N^{-1}\log N!$.  The active paper uses either a probability reference or,
when an absolute free energy is desired, an explicit ideal-gas
renormalisation.  The exact-number theorem also retains the cross covariance
between lattice and continuous constraints.

\subsection{Probability reference and finite pressure}

Let
\[
 \widetilde Z_{\varepsilon,N}(\lambda)
 ={1\over N!}\int_{\mathcal D_{\varepsilon,N}\times\mathbb R^{3N}}
 \exp\left\{\sum_{i=1}^N\lambda\cdot C(x_i,v_i)\right\}
 \,d\mathbf x\,d\mathbf v.                              \tag{B1.1}
\]
Fix an interior multiplier $\lambda_0$ and define the canonical probability
\[
 dP_{\varepsilon,N}^{0}
 ={1\over\widetilde Z_{\varepsilon,N}(\lambda_0)}
 {1\over N!}1_{\mathcal D_{\varepsilon,N}}
 \exp\left\{\sum_i\lambda_0\cdot C_i\right\}
 d\mathbf x\,d\mathbf v.                                \tag{B1.2}
\]
For a deterministic path/contact source $H$ and a centered constraint
multiplier $\theta$, put
\[
 Z_{\varepsilon,N}(H,\theta)
 =E_{P_{\varepsilon,N}^{0}}
 \exp\left\{\mu_\varepsilon H(\mathbf X)
 +\theta\cdot\left(\sum_iC_i-Nc_0\right)\right\},
 \quad q_{\varepsilon,N}={1\over N}\log Z_{\varepsilon,N}.             \tag{B1.3}
\]
Then $Z_{\varepsilon,N}(0,0)=1$ and
$q_{\varepsilon,N}(0,0)=0$ for every $N$.  If the absolute thermodynamic
free energy is required, we use
\[
 p_{\varepsilon,N}(\lambda)
 ={1\over N}\log\left(
 {N!\over |\mathbb T^3|^N}\widetilde Z_{\varepsilon,N}(\lambda)
 \right),                                                \tag{B1.4}
\]
so the Stirling term is removed explicitly rather than hidden in a theorem.

The grand-canonical generating functional is
\[
 \Xi_\varepsilon(z,H,\lambda)
 =\sum_{n\ge0}{z^n\over n!}
 \int_{\mathcal D_{\varepsilon,n}}
 e^{\sum_i\lambda\cdot C_i+\mu_\varepsilon H}
 \,d\mathbf x\,d\mathbf v.                              \tag{B1.5}
\]
B2-GC supplies its source-uniform analytic pressure before this paper is used.
Exact-$N$ quantities are coefficient ratios of (B1.5); the same $N!$ appears
in numerator and denominator and cannot create a $-\log N$ pressure.

\begin{proposition}[Normalization regression]
\label{prop:r23-b1-normalization}
On every compact source/multiplier set on which B2's pressure is finite,
$q_{\varepsilon,N}$ is locally bounded and any subsequential limit satisfies
$q(0,0)=0$.  The unnormalised pressure from (B1.1) is never called $q$.
\end{proposition}

\begin{proof}
The identity at the origin follows from (B1.2)--(B1.3).  Hölder's inequality
and the source exponential moment bound from B2 give local upper and lower
bounds.  The coefficient representation (B1.5) preserves the reference
normalisation.  In contrast, applying Stirling directly to (B1.1) would give
$-\log N+O(1)$; formula (B1.4) displays and removes precisely that term.
\end{proof}

\subsection{Quotient constraints and deterministic block exposure}

Let $G(\mathbf X)$ be the joint vector consisting of the exact particle
number/lattice constraints and the continuous momentum-energy/source
coordinates.  Quotient its target by all affine identities holding on
$\mathcal D_{\varepsilon,N}$ and call the resulting space $E=E_{\mathbb Z}
\oplus E_{\mathbb R}$.  Nondegeneracy below is asserted only on $E$.

Lift the symmetric integral to labelled particles and partition the labels,
not the realised configuration, into deterministic blocks
\[
 B_j=\{m(j-1)+1,\ldots,mj\},\qquad j=1,\ldots,\lfloor N/m\rfloor.         \tag{B1.6}
\]
The factor $1/N!$ is restored after the calculation by symmetry.  For each
block use a fixed finite atlas made from spatial grid boxes, velocity balls,
and full-rank minors of the block constraint map.  Its partition functions
are $C^s$, supported strictly inside their charts, and vanish with all
derivatives through order $s-1$ at chart boundaries.  The chart index is
expanded by the partition of unity in lexicographic block order.  Therefore
at step $j$ it is measurable with respect to the previously exposed blocks
and the current block; it is not selected after seeing the full
configuration.

\begin{lemma}[Boundary-flat one-block density]
\label{lem:r23-b1-block}
For a fixed block size $m=m(\dim E)$ and derivative order
$s>\dim E+2$, every conditional block amplitude is a finite sum of terms for
which a set of position--velocity coordinates $y\in\mathbb R^{\dim E}$ has
Jacobian determinant bounded below.  The push-forward of each term by the
quotient constraint map has a zero-extended density in $W^{s,1}(E_{\mathbb R})$
with a uniform norm.  The constants are uniform under all earlier block
conditionings in the regular canonical parameter compact.
\end{lemma}

\begin{proof}
The quotient removes exactly the directions in which every minor vanishes.
Compactness of the normalized constraint sphere gives a finite family of
nonzero minors.  Spatial grid refinement provides separated coordinate boxes
on which hard-core indicators involving the current block are constant in the
chosen interior variables; near a contact boundary an incoming normal
coordinate is used instead.  Coarea in the full position--velocity variables
gives the density.  Because each partition function and its first $s-1$
derivatives vanish at the chart boundary, extension by zero creates no
boundary distributions.  Repeated integration by parts therefore yields a
uniform $W^{s,1}$ norm.  Earlier blocks enter only as parameters in a compact
regular set; the finite atlas and lower minor bound are uniform there.
\end{proof}

A fixed positive fraction of the deterministic label blocks is usable after
conditioning: if too many blocks occupied only boundary collars or velocity
tails, the canonical energy and hard-core packing bounds would be violated.
The remaining blocks are covered by alternative charts rather than placed in
a frequency-independent exceptional event.

\subsection{Predictable product damping}

Let $\varphi_{j}(\xi)$ be the conditional characteristic function of block
$j$ in a chosen chart, where
$\xi=(u,v)\in[-\pi,\pi]^{\dim E_{\mathbb Z}}
\times\mathbb R^{\dim E_{\mathbb R}}$.  Lemma
\ref{lem:r23-b1-block} and covariance positivity on the quotient give
\[
 |\varphi_j(\xi)|\le
 \begin{cases}
 1-c|\xi|^2,&|\xi|\le\delta,\\
 1-c,&\delta\le|\xi|\le R,\\
 C(1+|v|)^{-s},&|v|>R.
 \end{cases}                                             \tag{B1.7}
\]
The block order (B1.6) is deterministic, and the chart partition is adapted to
the successive conditioning sigma-fields.  Multiplication of (B1.7) is
therefore a valid tower-property calculation.  Summing the finitely many chart
words preserves the same bounds because their nonnegative partition weights
sum to one.  We obtain the global majorant
\[
 |\Phi_N(\xi)|\le
 C\exp(-cN\min\{|\xi|^2,1\})
 +C_N(1+|v|)^{-s\theta N},                                \tag{B1.8}
\]
for some $\theta>0$.  There is no additive $e^{-cN}$ constant independent of
frequency over the unbounded domain.

\subsection{The full mixed lattice--continuous local theorem}

Under the tilted canonical law write
\[
 \Sigma=
 \begin{pmatrix}
 \Sigma_{\mathbb Z\mathbb Z}&\Sigma_{\mathbb Z\mathbb R}\\
 \Sigma_{\mathbb R\mathbb Z}&\Sigma_{\mathbb R\mathbb R}
 \end{pmatrix}                                           \tag{B1.9}
\]
for the covariance on $E$.  It is positive definite there.  Let
\[
 \Sigma_{\mathbb R\mid\mathbb Z}
 =\Sigma_{\mathbb R\mathbb R}
 -\Sigma_{\mathbb R\mathbb Z}
  \Sigma_{\mathbb Z\mathbb Z}^{-1}
  \Sigma_{\mathbb Z\mathbb R}.                          \tag{B1.10}
\]

\begin{theorem}[Source-uniform mixed local coefficient]
\label{thm:r23-b1-llt}
Let $(z_N,y_N)$ be central deviations in
$E_{\mathbb Z}\oplus E_{\mathbb R}$ of size $O(\sqrt N)$.  Uniformly for
sources in a smaller analytic ball,
\[
 P\{G_{\mathbb Z}=Nm_{\mathbb Z}+z_N,
       G_{\mathbb R}\in Nm_{\mathbb R}+dy_N\}
 =\frac{
 e^{-\frac1{2N}(z_N,y_N)^\mathsf T\Sigma^{-1}(z_N,y_N)}}
 {(2\pi N)^{\dim E/2}\sqrt{\det\Sigma}}
 (1+o(1))\,dy_N.                                         \tag{B1.11}
\]
Conditionally on the lattice value, the continuous deviation has mean
\[
 \Sigma_{\mathbb R\mathbb Z}
 \Sigma_{\mathbb Z\mathbb Z}^{-1}z_N                    \tag{B1.12}
\]
and covariance $N\Sigma_{\mathbb R\mid\mathbb Z}$.
\end{theorem}

\begin{proof}
Fourier invert on the torus in the lattice variables and on Euclidean space in
the continuous variables.  In the $N^{-1/2}$ ball, B2's analytic pressure has
the quadratic expansion with the full covariance (B1.9); rescaling yields the
joint Gaussian in (B1.11).  The compact annulus and unbounded continuous tail
are integrable by (B1.8).  Source derivatives are controlled by Cauchy bounds
on the same smaller ball.  Completing the square in the joint Gaussian gives
(B1.12) and (B1.10).  An unshifted marginal Gaussian would be correct only if
$\Sigma_{\mathbb R\mathbb Z}=0$, which is not assumed.
\end{proof}

\subsection{Activity saddle and exact-number extraction}

Let $P(z,H,\lambda)=\lim\mu_\varepsilon^{-1}\log
\Xi_\varepsilon(z,H,\lambda)$ be the B2-GC pressure.  The saddle
$(z_*,\lambda_*)$ solves the joint number/constraint equations in the quotient
space.  Strict positivity of (B1.9) gives a locally unique analytic saddle.
Use a Cauchy contour through $z_*$ and Fourier contours through
$\lambda_*$.  Theorem~\ref{thm:r23-b1-llt} controls the central chart, while
(B1.8) controls its complement.

\begin{theorem}[Exact-$N$ source ratio]
\label{thm:r23-b1-extraction}
For every regular central exact-number and mixed constraint window,
\[
 { [z^N]\Xi_\varepsilon(z,H,\lambda;\,G\in W_N)
  \over
   [z^N]\Xi_\varepsilon(z,0,\lambda_0;\,G\in W_N)}
 =\exp\{N(q(H)-q(0))+o(N)\}
 {\mathfrak g_H(W_N)\over\mathfrak g_0(W_N)}(1+o(1)),    \tag{B1.13}
\]
where $\mathfrak g_H$ is the mixed Gaussian with the conditional shift and
Schur covariance in (B1.10)--(B1.12).  Relative error is asserted only when
the Gaussian mass of $W_N$ dominates the absolute Fourier remainder.
\end{theorem}

\begin{proof}
Coefficient extraction from (B1.5) contributes the same $N!$ convention to
both numerator and denominator.  Deform the activity contour to the common
steepest-descent saddle.  Apply Theorem~\ref{thm:r23-b1-llt} jointly in the
activity, lattice, and continuous coordinates.  Outside the central chart,
(B1.8) is integrable and exponentially smaller.  Dividing the two expansions
gives (B1.13).  If a continuous window shrinks below the absolute inversion
error, the theorem states only the absolute estimate and does not manufacture
a relative one.
\end{proof}

\subsection{Round-Twenty-Two regression and literature position}

At $H=\theta=0$, the active canonical partition is exactly one, so the
$-\log N$ contradiction cannot occur.  Deterministic label blocks and adapted
finite atlases make conditional characteristic contractions genuinely
multiplicative.  Boundary-flat partitions prove, rather than assume,
zero-extended $W^{s,1}$ push-forward regularity.  Finally, the mixed local
formula uses the full covariance, conditional mean shift, and Schur
complement.

The grand-canonical pressure is imported only from the earlier B2-GC node, and
this paper returns exact-number conditioning to B2-MC.  The distinction between
absolute free energy and probability-normalized ensemble transfer is consistent
with standard ensemble theory \cite{Touchette2015}; the source-uniform
hard-sphere input is the connected history pressure proved in B2 rather than an
independence assumption.
